
\documentclass{article}
\usepackage{colortbl}
\usepackage{makecell}
\usepackage{multirow}
\usepackage{supertabular}

\begin{document}

\newcounter{utterance}

\centering \large Interaction Transcript for game `cladder', experiment `full\_v1.5\_default', episode 4386 with refineonly{-}Qwen3.5{-}9B.
\vspace{24pt}

{ \footnotesize  \setcounter{utterance}{1}
\setlength{\tabcolsep}{0pt}
\begin{supertabular}{c@{$\;$}|p{.15\linewidth}@{}p{.15\linewidth}p{.15\linewidth}p{.15\linewidth}p{.15\linewidth}p{.15\linewidth}}
    \# & \multicolumn{2}{c}{Player} && \multicolumn{2}{c}{Game Master} \\
    \hline

    \theutterance \stepcounter{utterance}  
    & & & \multicolumn{4}{p{0.6\linewidth}}{
        \cellcolor[rgb]{0.9,0.9,0.9}{
            \makecell[{{p{\linewidth}}}]{
                \texttt{\tiny{[P1$\langle$GM]}}
                \texttt{You are an expert in causal inference. The following question is not a typical commonsense query, but rather a meticulously designed question created by a professor specializing in causal inference, intended to assess the students' mastery of the course content.} \\
\texttt{Imagine a self{-}contained, hypothetical world with only the following conditions, and without any unmentioned factors or causal relationships: Gender has a direct effect on smoking and peanut allergy. Smoking has a direct effect on tar deposit. Tar deposit has a direct effect on peanut allergy. Gender is unobserved. For nonsmokers and with no tar deposit, the probability of being allergic to peanuts is 35\%. For nonsmokers and with high tar deposit, the probability of being allergic to peanuts is 63\%. For smokers and with no tar deposit, the probability of being allergic to peanuts is 25\%. For smokers and with high tar deposit, the probability of being allergic to peanuts is 49\%. For nonsmokers, the probability of high tar deposit is 22\%. For smokers, the probability of high tar deposit is 73\%. For smokers, would it be more likely to see being allergic to peanuts if the person had been a nonsmoker?} \\
\texttt{Start your answer with 'Yes' or 'No', followed by additional reasoning or evidence to support your explanation.} \\
            }
        }
    }
     \\ \\

    \theutterance \stepcounter{utterance}  
    & \multicolumn{4}{p{0.6\linewidth}}{
        \cellcolor[rgb]{0.561,0.737,0.561}{
            \makecell[{{p{\linewidth}}}]{
                \texttt{\tiny{[P1$\rangle$GM]}}
                \texttt{No} \\
\\ 
\texttt{**Reasoning:**} \\
\\ 
\texttt{To determine if a smoker would be more likely to be allergic to peanuts if they had been a nonsmoker, we must calculate the **Counterfactual Probability** \$P(\textbackslash{}text\{Allergy\} \textbackslash{}mid \textbackslash{}text\{Smoker\}, \textbackslash{}text\{do\}(\textbackslash{}text\{Nonsmoker\}))\$. This requires estimating the probability of having a peanut allergy for a smoker under the hypothetical intervention that they do not smoke.} \\
\\ 
\texttt{In this causal structure, the effect of smoking on peanut allergy is entirely mediated through **tar deposit**. There is no direct path from Smoking to Allergy that bypasses Tar Deposit. Therefore, the counterfactual probability for a smoker who did not smoke is equal to the probability of having a peanut allergy for a **nonsmoker** with the same distribution of tar deposits as a nonsmoker.} \\
\\ 
\texttt{Let's analyze the causal mechanism strictly based on the provided numbers:} \\
\texttt{1.  **Actual Scenario (Smoker):**} \\
\texttt{    *   \$P(\textbackslash{}text\{High Tar\} \textbackslash{}mid \textbackslash{}text\{Smoker\}) = 0.73\$} \\
\texttt{    *   \$P(\textbackslash{}text\{Low Tar\} \textbackslash{}mid \textbackslash{}text\{Smoker\}) = 0.27\$} \\
\texttt{    *   \$P(\textbackslash{}text\{Allergy\} \textbackslash{}mid \textbackslash{}text\{Smoker\}, \textbackslash{}text\{High Tar\}) = 0.49\$} \\
\texttt{    *   \$P(\textbackslash{}text\{Allergy\} \textbackslash{}mid \textbackslash{}text\{Smoker\}, \textbackslash{}text\{Low Tar\}) = 0.25\$} \\
\texttt{    *   Actual \$P(\textbackslash{}text\{Allergy\} \textbackslash{}mid \textbackslash{}text\{Smoker\}) = (0.73 \textbackslash{}times 0.49) + (0.27 \textbackslash{}times 0.25) = 0.3577 + 0.0675 = 0.4252\$ (42.52\%).} \\
\\ 
\texttt{2.  **Counterfactual Scenario (Smoker had been a Nonsmoker):**} \\
\texttt{    *   If a person stops smoking (becomes a nonsmoker), their probability of having high tar deposit changes. The prompt states: "For nonsmokers, the probability of high tar deposit is 22\%."} \\
\texttt{    *   Crucially, in a causal world where "Smoking has a direct effect on tar deposit," the counterfactual state of "Nonsmoker" implies the probability of high tar deposit reverts to the nonsmoker baseline (\$0.22\$), not the smoker baseline (\$0.73\$).} \\
\texttt{    *   We must use the allergy probabilities associated with the nonsmoking status for the specific tar levels. The prompt states: "For nonsmokers and with no tar deposit, the probability of being allergic to peanuts is 35\%." and "For nonsmokers and with high tar deposit,} \\
            }
        }
    }
    & & \\ \\

    \theutterance \stepcounter{utterance}  
    & & & \multicolumn{2}{p{0.3\linewidth}}{
        \cellcolor[rgb]{0.9,0.9,0.9}{
            \makecell[{{p{\linewidth}}}]{
                \texttt{\tiny{[GM$|$GM]}}
                \texttt{No} \\
            }
        }
    }
    & & \\ \\

    \theutterance \stepcounter{utterance}  
    & & & \multicolumn{2}{p{0.3\linewidth}}{
        \cellcolor[rgb]{0.9,0.9,0.9}{
            \makecell[{{p{\linewidth}}}]{
                \texttt{\tiny{[GM$|$GM]}}
                \texttt{no} \\
            }
        }
    }
    & & \\ \\

    \theutterance \stepcounter{utterance}  
    & & & \multicolumn{2}{p{0.3\linewidth}}{
        \cellcolor[rgb]{0.9,0.9,0.9}{
            \makecell[{{p{\linewidth}}}]{
                \texttt{\tiny{[GM$|$GM]}}
                \texttt{game\_result = WIN} \\
            }
        }
    }
    & & \\ \\

\end{supertabular}
}

\end{document}
